\chapter{Series de Laurent y teorema de los residuos}
\label{ch:b3:residues}

¿Qué le ocurre a una \omterm{def:b3:holomorphic:holo}{función holomorfa}
cerca de un punto donde \emph{no} está definida? La respuesta es una
tricotomía \omterm{def:b3:complete:complete}{completa} --- punto
evitable, polo o singularidad esencial --- que se lee en una serie de
potencias con dos colas, el desarrollo de Laurent. Un coeficiente de ese
desarrollo, el \emph{\omterm{def:b3:residues:singularities}{residuo}},
controla toda \omterm{def:b3:holomorphic:contour}{integral de contorno}
alrededor de la singularidad: el teorema de los
\omterm{def:b3:residues:singularities}{residuos} convierte integrales
definidas difíciles en álgebra finita, cuenta ceros de funciones
(principio del argumento, Rouché) y demuestra el teorema de la
aplicación abierta. Primero mejoramos el teorema de Cauchy, pasando de
los dominios estrellados a su forma definitiva, sin homología --- el
elegante argumento de Dixon ---, de modo que queden disponibles todos
los contornos de índice nulo respecto del complementario.

\section{El teorema global de Cauchy}

Un \emph{ciclo}\index{ciclo} $\Gamma$ es una suma formal finita de
\omterm{def:b3:topology:pathconnected}{caminos} cerrados
$\gamma_1, \dots, \gamma_m$; las integrales y los índices a lo largo de
$\Gamma$ son las sumas correspondientes, y
$\operatorname{im}\Gamma = \bigcup\operatorname{im}\gamma_j$.

\begin{theorem}[Cauchy, forma global]
\label{thm:b3:residues:globalcauchy}
Sean $\Omega \subseteq \C$ abierto, $f \in \mathcal H(\Omega)$ y
$\Gamma$ un ciclo en $\Omega$ tal que
\[
\operatorname{Ind}_\Gamma(w) = 0
\qquad\text{para todo } w \notin \Omega .
\]
Entonces, para todo $z \in \Omega\setminus\operatorname{im}\Gamma$,
\[
\frac1{2\iu\pi}\int_\Gamma\frac{f(w)}{w - z}\,\dd w =
\operatorname{Ind}_\Gamma(z)\,f(z),
\qquad\text{y}\qquad
\int_\Gamma f(w)\,\dd w = 0 .
\]
\index{teorema de Cauchy (global)}
\end{theorem}

\begin{proof}[Demostración (Dixon)]
Defínase $g \colon \Omega\times\Omega \to \C$ mediante
\[
g(z, w) = \begin{cases}
\dfrac{f(w) - f(z)}{w - z} & w \neq z,\\[4pt]
f'(z) & w = z .
\end{cases}
\]
\emph{$g$ es \omterm{def:b3:topology:continuity}{continua}}: fuera de la
diagonal, es claro. Cerca de un punto diagonal $(a, a)$, desarróllese
$f$ en serie de potencias en $a$
(el~\cref{thm:b3:holomorphic:analytic}): $f(w) - f(z) =
\sum_{n\geq1}c_n\bigl((w-a)^n - (z-a)^n\bigr)$, y dividiendo cada
término por $w - z$ (factorización de $u^n - v^n$) se obtiene, para
$z, w \in D(a, r)$,
\[
g(z, w) =
\sum_{n\geq1}c_n\sum_{j=0}^{n-1}(w-a)^{\,j}(z-a)^{\,n-1-j},
\]
válido también en la diagonal (cada suma \omterm{def:b3:topology:interior}{interior} se convierte en
$n(z-a)^{n-1}$, que suma $f'(z)$). Para $r$ pequeño, la serie converge
uniformemente en $D(a,r)^2$ ($\abs{\text{término}} \leq n\abs{c_n}r^{n-1}$,
sumable dentro del radio): la suma es
\omterm{def:b3:topology:continuity}{continua}.

Póngase $h(z) = \frac1{2\iu\pi}\int_\Gamma g(z, w)\,\dd w$ en $\Omega$:
es \omterm{def:b3:topology:continuity}{continua}
(\omterm{def:b3:topology:continuity}{continuidad} uniforme de $g$ en los
\omterm{def:b3:topology:compact}{compactos}) y
\omterm{def:b3:holomorphic:holo}{holomorfa} --- por Morera (el criterio
del~\cref{thm:b3:holomorphic:weierstrassconv}): para un triángulo
$T \subseteq \Omega$, Fubini da $\int_{\partial
T}h = \frac1{2\iu\pi}\int_\Gamma\bigl(\int_{\partial T}g(z,
w)\dd z\bigr)\dd w = 0$, y la integral \omterm{def:b3:topology:interior}{interior} se anula porque
$z \mapsto g(z, w)$ es \omterm{def:b3:holomorphic:holo}{holomorfa} en
$\Omega$ (en $z = w$ la singularidad es evitable: $g$ es
\omterm{def:b3:topology:continuity}{continua} allí y
\omterm{def:b3:holomorphic:holo}{holomorfa} en el resto --- el argumento
de extensión de la demostración del~\cref{thm:b3:holomorphic:formula}).

En el abierto
$\Omega' = \{z \notin \operatorname{im}\Gamma
: \operatorname{Ind}_\Gamma(z) = 0\}$, defínase $h_1(z) =
\frac1{2\iu\pi}\int_\Gamma\frac{f(w)}{w - z}\dd w$:
\omterm{def:b3:holomorphic:holo}{holomorfa} en $\Omega'$ (Morera o
derivación bajo la integral). Para $z \in \Omega\cap\Omega'$:
\[
h(z) = \frac1{2\iu\pi}\int_\Gamma\frac{f(w)}{w-z}\dd w
- f(z)\operatorname{Ind}_\Gamma(z) = h_1(z) .
\]
Por hipótesis, $\Omega\cup\Omega' = \C$ ($w \notin \Omega
\Rightarrow \operatorname{Ind}_\Gamma(w) = 0$), de modo que $h$ y $h_1$
se pegan en una función entera $H$. Como la componente no acotada del
complementario de $\operatorname{im}\Gamma$ está contenida en $\Omega'$
y $h_1(z) \to 0$ cuando $\abs z \to \infty$ (cota ML), $H$ es acotada y
tiende a $0$: Liouville (el~\cref{cor:b3:holomorphic:liouville}) da
$H \equiv 0$. Así, $h \equiv 0$ en $\Omega$, que es la fórmula integral.
Aplicándola, para $a \in
\Omega\setminus\operatorname{im}\Gamma$ fijo, a $\tilde f(w) =
(w - a)f(w)$ en $z = a$:
\[
\frac1{2\iu\pi}\int_\Gamma f(w)\dd w =
\frac1{2\iu\pi}\int_\Gamma \frac{\tilde f(w)}{w - a}\dd w =
\operatorname{Ind}_\Gamma(a)\,\tilde f(a) = 0 .
\]
\end{proof}

\section{Series de Laurent y singularidades aisladas}

\begin{theorem}[Desarrollo de Laurent]\label{thm:b3:residues:laurent}
Sea $f$ \omterm{def:b3:holomorphic:holo}{holomorfa} en la corona
$A = \{r < \abs{z - a} <
R\}$ ($0 \leq r < R \leq \infty$). Entonces\index{serie de Laurent}
\[
f(z) = \sum_{n\in\Z}c_n\,(z - a)^n
\qquad\text{en } A,
\qquad
c_n = \frac1{2\iu\pi}\int_{C_\rho}\frac{f(w)}{(w -
a)^{n+1}}\,\dd w
\]
para cualquier $r < \rho < R$ (con independencia de $\rho$), y las dos
semiseries convergen normalmente en las subcoronas
compactas. El desarrollo es único.
\end{theorem}

\begin{proof}
Fíjese $r < \rho_1 < \abs{z - a} < \rho_2 < R$ y sea $\Gamma =
C_{\rho_2} - C_{\rho_1}$ (la exterior en sentido positivo, la \omterm{def:b3:topology:interior}{interior}
en sentido negativo): un ciclo en $A$ con
$\operatorname{Ind}_\Gamma(w) = 0$ para todo $w \notin A$ (puntos
\omterm{def:b3:topology:interior}{interiores} al disco pequeño: $1 - 1 = 0$; exteriores al grande:
$0 - 0$). Por el~\cref{thm:b3:residues:globalcauchy},
$\operatorname{Ind}_\Gamma(z) = 1 - 0 = 1$ da
\[
f(z) = \frac1{2\iu\pi}\int_{C_{\rho_2}}\frac{f(w)}{w - z}\dd
w - \frac1{2\iu\pi}\int_{C_{\rho_1}}\frac{f(w)}{w - z}\dd w .
\]
Desarróllese el primer núcleo como en
el~\cref{thm:b3:holomorphic:analytic} (potencias de $\frac{z -
a}{w - a}$, de módulo $< 1$): la parte de índices no negativos
$\sum_{n\geq0}c_n(z-a)^n$. En el segundo, desarróllese al revés:
$\frac{-1}{w - z} = \frac1{(z-a)(1 - \frac{w - a}{z -
a})} = \sum_{m\geq0}\frac{(w-a)^m}{(z - a)^{m+1}}$, normalmente
convergente en $C_{\rho_1}$: la parte negativa
$\sum_{n\leq-1}c_n(z-a)^n$ con los coeficientes enunciados (índice
$n = -m-1$). Independencia de $\rho$: las integrales de los coeficientes
sobre $C_{\rho}$ y $C_{\rho'}$ difieren en $\int_\Gamma$ sobre un ciclo
de índice nulo de la \omterm{def:b3:holomorphic:holo}{holomorfa}
$\frac{f(w)}{(w-a)^{n+1}}$ en $A$: cero, de nuevo por
el~\cref{thm:b3:residues:globalcauchy}. Unicidad: intégrese
$\sum c_n(z-a)^n$ contra $(z - a)^{-m-1}$ sobre $C_\rho$ término a
término (convergencia normal): solo sobrevive $n = m$.
\end{proof}

\begin{definition}\label{def:b3:residues:singularities}
Si $f$ es \omterm{def:b3:holomorphic:holo}{holomorfa} en un disco
punteado $D(a, R)\setminus
\{a\}$, desarróllese por Laurent ($r = 0$). Tres casos
excluyentes:\index{singularidad aislada}
\begin{itemize}
  \item todos los $c_n = 0$ para $n < 0$:
  \emph{\omterm{thm:b3:residues:riemanncw}{singularidad evitable}} (la
  serie de índices no negativos extiende $f$
  \omterm{def:b3:holomorphic:holo}{holomorfamente} a $a$); \item
  $c_n \neq 0$ para un número finito, al menos uno, de $n < 0$: un
  \emph{polo} de orden $m = -\min\{n : c_n \neq 0\}$; equivalentemente,
  $f = g/(z-a)^m$, $g$ \omterm{def:b3:holomorphic:holo}{holomorfa},
  $g(a)
        \neq 0$; equivalentemente, $\abs{f(z)} \to \infty$ cuando $z\to
        a$; \item infinitos $c_n \neq 0$ negativos: \emph{singularidad
  esencial}.
\end{itemize}
El \emph{residuo}\index{residuo} es $\operatorname{Res}(f, a) = c_{-1}$.
Una \omterm{def:b3:holomorphic:holo}{función holomorfa} en $\Omega$
menos un conjunto de polos es \emph{meromorfa}\index{función meromorfa}
en $\Omega$.
\end{definition}

\begin{theorem}[Riemann; Casorati--Weierstrass]
\label{thm:b3:residues:riemanncw}
Sea $f$ \omterm{def:b3:holomorphic:holo}{holomorfa} en
$D(a,R)\setminus\{a\}$.
\begin{enumerate}
  \item (Riemann) Si $f$ está \emph{acotada} cerca de $a$, la
  singularidad es evitable.\index{singularidad evitable} \item
  (Casorati--Weierstrass) Si $a$ es esencial, entonces
  $f\bigl(D(a,\varepsilon)\setminus\{a\}\bigr)$ es denso en $\C$ para
  todo $\varepsilon$. \index{teorema de Casorati--Weierstrass}
\end{enumerate}
\end{theorem}

\begin{proof}
(1) Para $n < 0$ y $\rho \to 0$: $\abs{c_n} \leq
M\rho^{-n-1}\cdot\rho\cdot\rho^{-\,n}\dots$ por ML en $C_\rho$:
$\abs{c_n} \leq \frac{1}{2\pi}\,2\pi\rho\cdot
M\rho^{-(n+1)} = M\rho^{-n} \to 0$ (cuando $-n > 0$): todos los
coeficientes negativos se anulan. (2) Si algún valor $b$ no fuera
aproximado: $\abs{f - b} \geq
\delta$ cerca de $a$, de modo que $g = 1/(f - b)$ es
\omterm{def:b3:holomorphic:holo}{holomorfa} y acotada cerca de $a$:
evitable por (1), y $g$ se extiende con valor $c$. Si $c \neq 0$,
$f = b + 1/g$ está acotada cerca de $a$: evitable --- excluido. Si
$c = 0$, $g$ tiene un cero de orden finito $m$ en $a$
(el~\cref{thm:b3:holomorphic:identity}; $g \not\equiv 0$), y
$f = b + 1/g$ tiene un polo de orden $m$: excluido también.
\end{proof}

\section{El teorema de los residuos}

\begin{theorem}[Teorema de los residuos]\label{thm:b3:residues:residue}
Sean $\Omega$ abierto, $S \subseteq \Omega$ finito, $f \in
\mathcal H(\Omega\setminus S)$ y $\Gamma$ un ciclo en
$\Omega\setminus S$ con $\operatorname{Ind}_\Gamma(w) = 0$ para todo
$w \notin \Omega$. Entonces\index{teorema de los residuos}
\[
\frac{1}{2\iu\pi}\int_\Gamma f(z)\,\dd z
= \sum_{a\in S}\operatorname{Ind}_\Gamma(a)\,
\operatorname{Res}(f, a) .
\]
\end{theorem}

\begin{proof}
Para cada $a \in S$, sea $P_a(z) = \sum_{n\leq-1}c_n^{(a)}(z -
a)^n$ la parte principal de $f$ en $a$: una serie convergente en
$\C\setminus\{a\}$ (su radio en $1/(z-a)$ es infinito: la cola de
Laurent converge para todo $\abs{z-a}$ pequeño y, siendo una serie de
potencias en $(z-a)^{-1}$, converge en todas partes),
\omterm{def:b3:holomorphic:holo}{holomorfa} allí. Entonces
$g = f - \sum_{a\in S}P_a$ tiene \omterm{thm:b3:residues:riemanncw}{singularidades evitables} en cada punto
de $S$ (su desarrollo de Laurent en $a$ no tiene parte negativa: las
demás $P_{a'}$ son \omterm{def:b3:holomorphic:holo}{holomorfas} en $a$),
de modo que $g$ se extiende
\omterm{def:b3:holomorphic:holo}{holomorfamente} a $\Omega$, y
el~\cref{thm:b3:residues:globalcauchy} da $\int_\Gamma g = 0$. Queda
integrar cada $P_a$: término a término (convergencia normal en el
\omterm{def:b3:topology:compact}{compacto} $\operatorname{im}\Gamma$,
que evita $a$),
\[
\frac1{2\iu\pi}\int_\Gamma(z - a)^n\,\dd z = 0 \ (n \leq -2:
\text{primitiva } \tfrac{(z-a)^{n+1}}{n+1}),
\qquad
\frac1{2\iu\pi}\int_\Gamma\frac{\dd z}{z - a} =
\operatorname{Ind}_\Gamma(a),
\]
de modo que $\frac1{2\iu\pi}\int_\Gamma P_a =
c_{-1}^{(a)}\operatorname{Ind}_\Gamma(a)$. Súmese sobre $a$.
\end{proof}

\begin{method}[Cálculo de residuos]\label{met:b3:residues:compute}
Polo simple: $\operatorname{Res}(f, a) = \lim_{z\to a}(z -
a)f(z)$; para $f = g/h$ con $g(a) \neq 0$, $h(a) = 0$, $h'(a)
\neq 0$: $\operatorname{Res} = g(a)/h'(a)$. Polo de orden $m$:
$\operatorname{Res}(f, a) = \frac1{(m-1)!}\lim_{z\to
a}\bigl((z-a)^mf(z)\bigr)^{(m-1)}$. Singularidades esenciales:
desarróllese y léase $c_{-1}$ (por ejemplo, a partir de series
conocidas). Verifíquese siempre qué polos rodea realmente el contorno, y
con qué índice.
\end{method}

\begin{example}[Los cuatro tipos clásicos de integral]
\label{ex:b3:residues:integrals}
(a) \emph{Racional sobre $\R$}: para $\int_\R\frac{\dd x}{1 +
x^4}$, ciérrese con una semicircunferencia grande $S_R$ en el semiplano
superior: allí el integrando es $O(R^{-4})$, de modo que
$\int_{S_R} \to 0$ (ML), y el teorema de los
\omterm{def:b3:residues:singularities}{residuos} con los polos
$\eu^{\iu\pi/4}, \eu^{3\iu\pi/4}$ (simples,
\omterm{def:b3:residues:singularities}{residuos}
$\frac1{4z^3} = \frac{z}{4z^4} = -\frac z4$ en un polo) da
\[
\int_\R\frac{\dd x}{1 + x^4}
= 2\iu\pi\Bigl(-\frac{\eu^{\iu\pi/4}}4 -
\frac{\eu^{3\iu\pi/4}}4\Bigr)
= \frac{\pi}{\sqrt2} .
\]
(b) \emph{Tipo Fourier}: para $t \geq 0$,
$\int_\R\frac{\eu^{\iu tx}}{1 + x^2}\dd x =
2\iu\pi\operatorname{Res}\bigl(\tfrac{\eu^{\iu tz}}{1+z^2},
\iu\bigr) = 2\iu\pi\frac{\eu^{-t}}{2\iu} = \pi\eu^{-t}$ ---
la semicircunferencia superior funciona porque allí
$\abs{\eu^{\iu tz}} =
\eu^{-t\operatorname{Im}z} \leq 1$; tomando partes reales:
$\int_\R\frac{\cos(tx)}{1+x^2}\dd x = \pi\eu^{-\abs t}$,
lo que salda la fórmula admitida en el~\cref{exo:b3:lebesgue:10}. (c)
\emph{Trigonométrica sobre un periodo}: sustitúyase $z =
\eu^{\iu t}$, $\cos t = \frac{z + z^{-1}}2$, $\dd t =
\frac{\dd z}{\iu z}$: $\int_0^{2\pi}\frac{\dd t}{a + \cos t}$ ($a > 1$)
se convierte en un recuento de
\omterm{def:b3:residues:singularities}{residuos} dentro de la
circunferencia unidad
(\cref{exo:b3:residues:2}).
(d) \emph{Series}: emparéjese $f$ con $\pi\cot(\pi z)$, cuyos polos son
los enteros con \omterm{def:b3:residues:singularities}{residuo} $1$: el
problema de fin de semana suma así $\sum n^{-2}$ y $\sum n^{-4}$.
\end{example}

\begin{omfigure}
\begin{tikzpicture}[scale=1.15]
  \draw[->] (-2.5,0) -- (2.6,0) node[right, font=\small]
    {$\operatorname{Re}$};
  \draw[->] (0,-0.5) -- (0,2.5) node[above, font=\small]
    {$\operatorname{Im}$};
  \draw[omThm, very thick, ->] (-2,0) -- (0.3,0);
  \draw[omThm, very thick] (0.3,0) -- (2,0);
  \draw[omThm, very thick] (2,0) arc (0:90:2);
  \draw[omThm, very thick, ->] (2,0) arc (0:50:2);
  \draw[omThm, very thick] (0,2) arc (90:180:2);
  \node[omThm, font=\small] at (1.75,1.55) {$S_R$};
  \node[font=\small] at (2.2,-0.25) {$R$};
  \node[font=\small] at (-2.15,-0.25) {$-R$};
  \fill[omDef] (45:1.19) circle (1.6pt)
    node[above right, font=\small] {$\eu^{\iu\pi/4}$};
  \fill[omDef] (135:1.19) circle (1.6pt)
    node[above left, font=\small] {$\eu^{3\iu\pi/4}$};
  \fill[black!40] (-45:1.19) circle (1.6pt);
  \fill[black!40] (-135:1.19) circle (1.6pt);
\end{tikzpicture}

{\small El contorno semicircular para $\int_\R\frac{\dd x}{1 +
x^4}$: cuando $R \to \infty$, el arco aporta $O(R^{-3})$, y el teorema
de los \omterm{def:b3:residues:singularities}{residuos} cuenta los dos
polos encerrados (en azul). Los dos polos inferiores (en gris) quedan
fuera: índice $0$.}
\end{omfigure}

\section{El principio del argumento y el teorema de Rouché}

\begin{theorem}[Principio del argumento]
\label{thm:b3:residues:argument}
Sea $f$ \omterm{def:b3:residues:singularities}{meromorfa} en $\Omega$,
con ceros $z_j$ (de órdenes $m_j$) y polos $p_k$ (de órdenes $\mu_k$), y
sea $\gamma$ un \omterm{def:b3:topology:pathconnected}{camino} cerrado
en $\Omega$ que los evita todos, con $\operatorname{Ind}_\gamma = 0$
fuera de $\Omega$. Entonces \index{principio del argumento}
\[
\frac1{2\iu\pi}\int_\gamma\frac{f'(z)}{f(z)}\,\dd z
= \sum_j m_j\operatorname{Ind}_\gamma(z_j)
- \sum_k \mu_k\operatorname{Ind}_\gamma(p_k)
\]
(solo un número finito de términos es no nulo). Para un contorno simple
recorrido en sentido positivo, la integral cuenta los ceros menos los
polos \omterm{def:b3:topology:interior}{interiores}, con multiplicidad --- y es igual al índice del
\omterm{def:b3:topology:pathconnected}{camino} imagen $f\circ\gamma$
alrededor de $0$.
\end{theorem}

\begin{proof}
Cerca de un cero de orden $m$: $f = (z-a)^mg$, $g(a) \neq 0$, de modo
que $\frac{f'}f = \frac m{z - a} + \frac{g'}g$ con el segundo término
\omterm{def:b3:holomorphic:holo}{holomorfo} cerca de $a$: un polo simple
de \omterm{def:b3:residues:singularities}{residuo} $m$. Cerca de un polo
de orden $\mu$: $f = (z-a)^{-\mu}g$ da
\omterm{def:b3:residues:singularities}{residuo} $-\mu$. En el resto,
$\frac{f'}f$ es \omterm{def:b3:holomorphic:holo}{holomorfa}. (Los ceros
y polos de índice no nulo están en una región
compacta rodeada por $\gamma$; por el
principio de identidad son allí finitos en número, $f \not\equiv 0$.)
Aplíquese el~\cref{thm:b3:residues:residue}. La última observación:
$\frac1{2\iu\pi}\int_\gamma\frac{f'}f =
\frac1{2\iu\pi}\int_{f\circ\gamma}\frac{\dd w}w =
\operatorname{Ind}_{f\circ\gamma}(0)$ (sustitúyase $w =
f(\gamma(t))$).
\end{proof}

\begin{theorem}[Rouché]\label{thm:b3:residues:rouche}
Sean $f, g$ \omterm{def:b3:holomorphic:holo}{holomorfas} en $\Omega$ y
$\gamma$ un \omterm{def:b3:topology:pathconnected}{camino} cerrado con
$\operatorname{Ind}_\gamma \in \{0,1\}$, nulo fuera de $\Omega$ (un
contorno simple). Si\index{teorema de Rouché}
\[
\abs{g(z)} < \abs{f(z)} \qquad \text{en }
\operatorname{im}\gamma,
\]
entonces $f$ y $f + g$ tienen el mismo número de ceros (con
multiplicidad) en la región $\{\operatorname{Ind}_\gamma =
1\}$.
\end{theorem}

\begin{proof}
Para $t \in \intcc01$, $f_t = f + tg$ no tiene ningún cero en
$\operatorname{im}\gamma$ ($\abs{f_t} \geq \abs f - \abs g >
0$), de modo que
\[
N(t) = \frac1{2\iu\pi}\int_\gamma
\frac{f_t'(z)}{f_t(z)}\,\dd z
\]
está bien definida; cuenta los ceros en la región encerrada
(el~\cref{thm:b3:residues:argument}; sin polos). $N$ es
\omterm{def:b3:topology:continuity}{continua} en $t$ (el integrando es
\omterm{def:b3:topology:continuity}{continuo} conjuntamente y los
denominadores están uniformemente acotados inferiormente ---
convergencia dominada) y toma valores enteros: constante. $N(0) = N(1)$.
\end{proof}

\begin{corollary}[Teorema de la aplicación abierta]
\label{cor:b3:residues:openmapping}
Una \omterm{def:b3:holomorphic:holo}{función holomorfa} no constante en
un abierto
\omterm{def:b3:topology:connected}{conexo} es una aplicación abierta. En
particular (una vez más) vale el principio del máximo, y una biyección
\omterm{def:b3:holomorphic:holo}{holomorfa} tiene inversa
\omterm{def:b3:holomorphic:holo}{holomorfa}.\index{teorema de la
aplicación abierta (holomorfa)}
\end{corollary}

\begin{proof}
Sea $f(a) = b$; $f - b$ tiene un cero de cierto orden finito $m
\geq 1$ en $a$ (principio de identidad: $f \not\equiv b$). Elíjase $r$
con $f - b$ sin ceros en $\bar D(a, r)\setminus\{a\}$
(\omterm{thm:b3:holomorphic:identity}{ceros aislados}) y sea
$\delta = \min_{\abs{z - a} =
r}\abs{f(z) - b} > 0$. Para $\abs{w - b} < \delta$: en la
circunferencia, $\abs{(b - w)} < \delta \leq \abs{f - b}$, de modo que
Rouché ($f - b$ frente a la constante $b - w$) dice que $f - w$ tiene
exactamente $m$ ceros en $D(a, r)$: todo tal $w$ se alcanza ---
$f(D(a,r)) \supseteq D(b, \delta)$: abierta. Principio del máximo: un
máximo \omterm{def:b3:topology:interior}{interior} de $\abs f$ es
imposible para $f$ no constante, pues su imagen alrededor de $f(a)$
contiene puntos de módulo mayor. Inversa: una biyección
\omterm{def:b3:holomorphic:holo}{holomorfa} $f$ es abierta, de modo que
$f^{-1}$ es \omterm{def:b3:topology:continuity}{continua}; el cero de
$f - f(a)$ en $a$ es simple ($m \geq 2$ daría $m$ preimágenes de los
valores próximos --- distintas, pues $f'$ solo se anula en puntos
aislados, de modo que cerca de $a$ los $m$ ceros de $f - w$ son simples
y distintos para $w$ pequeño genérico: contradicción con la
inyectividad); entonces $f'(a) \neq 0$ y el cociente incremental de
$f^{-1}$ converge: $\bigl(f^{-1}\bigr)'(b) =
1/f'(a)$.
\end{proof}

\section{Ejercicios}

\begin{exercise}[$\star$]\label{exo:b3:residues:1}
Clasifíquese la singularidad en $0$ y calcúlese el
\omterm{def:b3:residues:singularities}{residuo}:
\[
\frac{\sin z}{z},\qquad
\frac{\eu^z - 1}{z^2},\qquad
\frac{1}{z(z-1)^2},\qquad
\frac{\cos z}{z^3},\qquad
\eu^{1/z},\qquad
\frac1{\sin z} .
\]
Dese además el \omterm{def:b3:residues:singularities}{residuo} de la
tercera en $z = 1$ y el de la última en $z = \pi$.
\end{exercise}

\begin{exercise}[$\star$]\label{exo:b3:residues:2}
Para $a > 1$ calcúlese, vía $z = \eu^{\iu t}$:
\[
\int_0^{2\pi}\frac{\dd t}{a + \cos t}
= \frac{2\pi}{\sqrt{a^2 - 1}} .
\]
Compruébense los comportamientos límite $a \to 1^+$ y $a \to \infty$.
\end{exercise}

\begin{exercise}[$\star\star$]\label{exo:b3:residues:3}
Calcúlense con contornos semicirculares, justificando las estimaciones
sobre los arcos:
\[
\int_\R\frac{x^2}{1 + x^6}\,\dd x = \frac\pi3,
\qquad
\int_\R\frac{\dd x}{(1 + x^2)^{2}} = \frac\pi2
\quad\text{(recuperando el \cref{exo:b3:fouriertransform:3})}.
\]
\end{exercise}

\begin{exercise}[$\star\star$]\label{exo:b3:residues:4}
Demuéstrese, para $t \geq 0$ y $a > 0$:
\[
\int_\R\frac{\cos(tx)}{x^2 + a^2}\,\dd x =
\frac{\pi}{a}\,\eu^{-at},
\]
y dedúzcase por inversión la transformada de Fourier de
$x \mapsto \eu^{-a\abs
x}$ --- comparando con
\cref{exo:b3:fouriertransform:1}.
\end{exercise}

\begin{exercise}[$\star\star$]\label{exo:b3:residues:5}
Desarróllese $f(z) = \dfrac1{(z-1)(z-2)}$ en
\omterm{thm:b3:residues:laurent}{serie de Laurent} en cada una de las
tres regiones $\abs z < 1$, $1 < \abs z < 2$, $\abs z >
2$. ¿Por qué difieren los tres desarrollos? Explíquese por qué el
coeficiente de $z^{-1}$ en el segundo y el tercer desarrollos \emph{no}
es un \omterm{def:b3:residues:singularities}{residuo} de $f$ en $0$ ($f$
no tiene allí ninguna singularidad), y calcúlense los
\omterm{def:b3:residues:singularities}{residuos} verdaderos de $f$, en
$1$ y en
$2$.
\end{exercise}

\begin{exercise}[$\star\star$]\label{exo:b3:residues:6}
(a) Demuéstrese que $\eu^{1/z}$ tiene una singularidad esencial en $0$ y
verifíquese Casorati--Weierstrass a mano: resuélvase $\eu^{1/z} =
w$ explícitamente para cualquier $w \neq 0$, exhibiendo soluciones
arbitrariamente próximas a $0$. (b) Demuéstrese que $\abs{\eu^{1/z}}$ no
está acotada en ningún \omterm{def:b3:topology:topology}{entorno}
punteado de $0$ y, sin embargo, $\eu^{1/z}$ no tiene ningún polo: ¿qué
límite falla?
\end{exercise}

\begin{exercise}[$\star\star$]\label{exo:b3:residues:7}
Cuéntense con Rouché: (a) los ceros de $z^7 - 4z^3 + z - 1$ en
$\abs z < 1$; (b) los ceros de $z^4 + 5z + 1$ en $\abs z < 1$ y en $1 <
\abs z < 2$; (c) demuéstrese de nuevo d'Alembert--Gauss: un polinomio
mónico de grado $n$ tiene $n$ ceros en cierto disco grande
\emph{(compárese con
$z^n$)}.
\end{exercise}

\begin{exercise}[$\star\star\star$]\label{exo:b3:residues:8}
(Hurwitz) Sea $f_n \to f$ uniformemente en los
\omterm{def:b3:topology:compact}{compactos}, con $f_n \in
\mathcal H(\Omega)$, $\Omega$
\omterm{def:b3:topology:connected}{conexo}, $f \not\equiv 0$. (a)
Demuéstrese que si todas las $f_n$ carecen de ceros, también $f$.
\emph{(Si $f(a) = 0$: principio del argumento en una circunferencia
pequeña alrededor de $a$, y el~\cref{thm:b3:holomorphic:weierstrassconv}
para pasar al límite en $\int f_n'/f_n$.)} (b) Demuéstrese que si todas
las $f_n$ son inyectivas, entonces $f$ es inyectiva o constante.
\emph{(Aplíquese (a) a $z \mapsto f_n(z) - f_n(w)$ en
$\Omega\setminus\{w\}$.)}
\end{exercise}

\begin{exercise}[$\star\star\star$]\label{exo:b3:residues:9}
Para $n \geq 2$, intégrese $\frac1{1 + z^n}$ sobre la frontera del
sector $\{0 \leq \arg z \leq \frac{2\pi}n,\ \abs z
\leq R\}$ y dedúzcase
\[
\int_0^{+\infty}\frac{\dd x}{1 + x^n} =
\frac{\pi}{n\,\sin(\pi/n)} .
\]
Verifíquese $n = 2$ frente a $\arctan$, y el límite $n \to
\infty$.
\end{exercise}

\begin{exercise}[$\star\star$]\label{exo:b3:residues:10}
Sea $f$ una función racional con $\deg(\text{denominador})
\geq \deg(\text{numerador}) + 2$. Demuéstrese que la suma de
\emph{todos} los \omterm{def:b3:residues:singularities}{residuos} de $f$
es nula \emph{(intégrese sobre circunferencias cada vez mayores)}. Úsese
esto para recalcular la descomposición en fracciones simples de
$\frac1{z(z-1)(z-2)}$ sin álgebra lineal.
\end{exercise}

\begin{exercise}[$\star\star\star$]\label{exo:b3:residues:11}
(El ojo de cerradura: la integral de reflexión de Euler) Para
$0 < a < 1$, calcúlese
\[
I(a) = \int_0^{\infty}\frac{x^{a-1}}{1 + x}\,\dd x =
\frac{\pi}{\sin(\pi a)}
\]
integrando $f(z) = \frac{z^{a-1}}{1+z} =
\frac{\eu^{(a-1)\log z}}{1 + z}$ (logaritmo con corte a lo largo de
$\R_+$, $\arg z \in \intoo0{2\pi}$) sobre el contorno de ojo de
cerradura: hacia fuera por encima del corte, de $\varepsilon$ a $R$,
alrededor de $C_R$, de vuelta por debajo del corte y alrededor de
$C_\varepsilon$. Justifíquese: los dos tramos rectos difieren en el
factor $\eu^{2\iu\pi(a-1)}$, las aportaciones de las circunferencias se
anulan ($R^{a-1}\cdot R \to 0$ y
$\varepsilon^{a-1}\cdot\varepsilon \to 0$) y el único polo $z = -1$
tiene \omterm{def:b3:residues:singularities}{residuo}
$\eu^{\iu\pi(a - 1)}$. Dedúzcase también
$\Gamma(a)\Gamma(1 - a) = \frac\pi{\sin\pi a}$ \emph{(escríbase
$\Gamma(a)\Gamma(1-a) = B(a, 1-a)$ por el~\cref{pb:b3:lebesgue:1} y
sustitúyase $t =
\frac{x}{1+x}$)}.
\end{exercise}

\begin{exercise}[$\star\star$]\label{exo:b3:residues:12}
(Recuento de ceros con el principio del argumento, numéricamente) Sea
$P(z) = z^4 + 8z + 1$. (a) ¿Cuántos ceros hay en el disco unidad?
\emph{(Rouché frente a
$8z + 1$.)}
(b) ¿Cuántos en la corona $1 < \abs z < 3$? \emph{(Rouché frente a $z^4$
en $\abs z = 3$.)} Afínese: demuéstrese que todo cero tiene módulo
$< 2.1$. (c) ¿Cuántos en el semiplano derecho? \emph{(Cuéntense primero
en $\abs z = 2$; síganse después las imágenes del eje imaginario:
$P(\iu t) = t^4 + 1 + 8\iu t$ tiene parte real positiva en todo él, de
modo que no hay ceros sobre el eje, y la variación del argumento a lo
largo de él es calculable --- conclúyase con un semidisco grande.)}
\end{exercise}

\section{Problema: \texorpdfstring{$\zeta(2k)$}{zeta(2k)} mediante la
cotangente}

\begin{problem}[{Problema de fin de semana --- sumar $\sum n^{-2k}$
con
\omterm{def:b3:residues:singularities}{residuos}}]\label{pb:b3:residues:1}
El teorema de los \omterm{def:b3:residues:singularities}{residuos} suma
series: emparejar una función racional con $\pi\cot(\pi z)$, cuyos polos
están en los enteros, convierte $\sum_{n}f(n)$ en un recuento de
\omterm{def:b3:residues:singularities}{residuos}. Demostraremos el
método y calcularemos $\zeta(2) = \frac{\pi^2}{6}$ y $\zeta(4) =
\frac{\pi^4}{90}$ --- los valores hallados por series de Fourier en
segundo año y por trazas de operadores en el~\cref{ch:b3:spectral},
ahora por integración de contorno.

\medskip
\noindent\textbf{Parte I --- El núcleo cotangente.}
\begin{enumerate}
  \item Demostrar que $\pi\cot(\pi z)$ es
  \omterm{def:b3:residues:singularities}{meromorfa} en $\C$ con polos
  simples exactamente en $z = n \in \Z$, cada uno de
  \omterm{def:b3:residues:singularities}{residuo} $1$ \emph{(calcúlese
        $\lim_{z\to n}(z-n)\pi\cot\pi z$)}.
  \item Calcular el principio del desarrollo de Laurent en $0$:
        \[
        \pi\cot(\pi z) = \frac1z - \frac{\pi^2}{3}\,z -
        \frac{\pi^4}{45}\,z^3 + O(z^5) ,
        \]
        dividiendo la serie de potencias de $\cos$ por la de $\sin$
        (justifíquese la división: $\frac{\sin\pi z}{\pi z}$ es
        \omterm{def:b3:holomorphic:holo}{holomorfa} y no nula cerca de
        $0$, de modo que su inversa es
        \omterm{def:b3:holomorphic:holo}{holomorfa}; identifíquense los
        coeficientes hasta el orden $3$). \item Sea $C_N$ la frontera
        del cuadrado de vértices $(\pm1\pm\iu)(N + \frac12)$. Demostrar
        que $\abs{\cot(\pi z)} \leq 2$ en $C_N$ para todo $N \geq
        1$. \emph{(En los lados verticales, $\cot(\pi(\pm(N +
        \frac12) + \iu y)) = \mp\tan(\iu\pi y)$, de módulo
        $\abs{\tanh(\pi y)} \leq 1$; en los horizontales
        $\abs y = N + \frac12$, acótese
        $\abs{\cot(\pi(x\pm\iu y))} \leq \coth(\pi y) \leq
        \coth(\pi/2) < 1.1$.)}
\end{enumerate}

\medskip
\noindent\textbf{Parte II --- El teorema de sumación.}
\begin{enumerate}[resume]
  \item Sea $f$ racional, \omterm{def:b3:holomorphic:holo}{holomorfa} en
  los enteros, con $\deg(\text{denom}) \geq \deg(\text{num}) + 2$.
  Usando el teorema de los
  \omterm{def:b3:residues:singularities}{residuos} en $C_N$ y la cota de
  la pregunta 3, demuéstrese:
        \[
        \lim_{N\to\infty}\ \sum_{n = -N}^{N} f(n)
        = -\sum_{p\ \text{polo de}\ f}
        \operatorname{Res}\bigl(\pi\cot(\pi z)f(z),\,p\bigr).
        \]
  \item ¿Dónde necesita el argumento la condición sobre los grados?
  Muéstrese con un ejemplo (tómese $f(z) = 1/(z + \frac12)$) que, con
  decaimiento más lento, el límite simétrico puede existir aunque la
  serie bilátera diverja --- y que la fórmula calcula entonces el
  \emph{valor principal}.
\end{enumerate}

\medskip
\noindent\textbf{Parte III --- Los valores.}
\begin{enumerate}[resume]
  \item Aplíquese el método a $f(z) = 1/z^2$: aquí $f$ tiene su polo
  \emph{en} un entero, de modo que hágase el argumento directamente ---
  intégrese $g(z) = \frac{\pi\cot(\pi
        z)}{z^2}$ sobre $C_N$, demuéstrese que la integral $\to 0$ y
  calcúlese $\operatorname{Res}(g, 0)$ a partir de la pregunta 2.
  Conclúyase:
        \[
        2\sum_{n\geq1}\frac1{n^2} = \frac{\pi^2}{3},
        \qquad \zeta(2) = \frac{\pi^2}6 .
        \]
  \item Lo mismo con $g(z) = \frac{\pi\cot(\pi z)}{z^4}$: calcúlese
  $\operatorname{Res}(g, 0)$ y dedúzcase
        $\zeta(4) = \frac{\pi^4}{90}$.
  \item Explíquese el patrón general: para todo $k \geq 1$, $\zeta(2k)$
  es $-\frac12$ veces el coeficiente de $z^{2k-1}$ en el desarrollo de
  Laurent de $\pi\cot(\pi
        z)$ en $0$ --- un múltiplo racional de $\pi^{2k}$. Calcúlese
  $\zeta(6)$ llevando la división de la pregunta 2 un paso más allá.
  ¿Qué dice el método sobre $\zeta(3)$ --- y por qué no dice nada?
\end{enumerate}

\medskip
\noindent\textbf{Parte IV --- El desarrollo en fracciones simples de la
cotangente.}
\begin{enumerate}[resume]
  \item Fíjese $w \in \C\setminus\Z$ y aplíquese el método de la Parte
  II a $f(z) = \dfrac{1}{(z - w)(z + w)}$ --- obsérvese que
  $\pi\cot(\pi z)f(z)$ tiene ahora polos simples adicionales en $\pm w$,
  cuyos \omterm{def:b3:residues:singularities}{residuos} han de
  incorporarse al recuento. Dedúzcase el desarrollo en fracciones
  simples
        \[
        \pi\cot(\pi w) = \frac1w +
        \sum_{n\geq1}\frac{2w}{w^2 - n^2},
        \]
        convergiendo la serie normalmente en los
        \omterm{def:b3:topology:compact}{compactos} de
        $\C\setminus\Z$.
  \item Recupérense de este desarrollo, desarrollando cada término en
  potencias de $w$ (justifíquese el intercambio), los \emph{mismos}
  coeficientes de Laurent de la pregunta 2 --- el círculo se cierra: la
  fórmula de Euler $\sum\frac1{n^2} = \frac{\pi^2}6$ es el coeficiente
  de $w$ en las dos caras de la cotangente. Compárese con la
  demostración por series de Fourier (segundo año) y con la demostración
  por trazas (el~\cref{pb:b3:spectral:1}): tres teorías, un número.
\end{enumerate}

\medskip
\noindent\textbf{Parte V --- El producto de Euler para el seno.} El
desarrollo de la pregunta 9 es la derivada logarítmica de un producto
infinito; demostremos ahora honestamente la factorización de Euler de
1734.
\begin{enumerate}[resume]
  \item Para $N \geq 1$, póngase $P_N(z) =
        z\prod_{n=1}^{N}\bigl(1 - \frac{z^2}{n^2}\bigr)$. Demuéstrese
  que $P_N$ converge, uniformemente en cada disco $\bar D(0, R)$, a una
  función entera $P$ cuyos ceros son exactamente los enteros, todos
  simples. \emph{(Para $n
        \geq 2R$, escríbase el factor como $\exp\log(1 -
        z^2/n^2)$ con el logaritmo principal
  del~\cref{exo:b3:holomorphic:3}, acótese $\abs{\log(1+u)}
        \leq 2\abs u$ para $\abs u \leq \frac12$ mediante la serie y
  exponénciese la suma normalmente convergente de logaritmos; los
  factores restantes, en número finito, forman un polinomio. Conclúyase
  con
        \cref{thm:b3:holomorphic:weierstrassconv}.)}
  \item Demuéstrese que en $\C\setminus\Z$,
        \[
        \frac{P'(z)}{P(z)} = \frac1z +
        \sum_{n\geq1}\frac{2z}{z^2 - n^2} = \pi\cot(\pi z)
        \]
        \emph{(derívense los productos finitos, pásese al límite usando
        el~\cref{thm:b3:holomorphic:weierstrassconv} y la ausencia de
        ceros de $P$ fuera de $\Z$, y cítese la pregunta
        9)}.
  \item Demuéstrese que $Q = \sin(\pi z)/P(z)$ se extiende a una función
  entera sin ceros con $Q' = 0$, y conclúyase el \emph{producto de
  Euler}:
        \[
        \sin(\pi z) = \pi z\prod_{n\geq1}
        \Bigl(1 - \frac{z^2}{n^2}\Bigr)
        \qquad (z \in \C) .
        \]
  \item (Wallis, 1655) Evalúese en $z = \frac12$:
        \[
        \frac\pi2 = \prod_{n\geq1}\frac{4n^2}{4n^2 - 1}
        = \lim_{N\to\infty}
        \frac{2\cdot2\cdot4\cdot4\cdots(2N)(2N)}
        {1\cdot3\cdot3\cdot5\cdots(2N-1)(2N+1)} .
        \]
  \item Para $\abs z < 1$, desarróllese el logaritmo del producto como
  serie doble (justifíquese la reordenación) y recupérese $\zeta(2) =
        \frac{\pi^2}6$ identificando el coeficiente de $z^3$ en
  $\sin(\pi z) = \pi z - \frac{\pi^3}6z^3 + \cdots$ --- la cara
  multiplicativa del número de Euler.
\end{enumerate}

\medskip
\noindent\textbf{Parte VI --- Núcleos hermanos.} La cotangente tiene
parientes; cada uno pone precio a su propia familia de series.
\begin{enumerate}[resume]
  \item Derívese término a término el desarrollo de la pregunta 9
  (justificado por el~\cref{thm:b3:holomorphic:weierstrassconv}) para
  obtener, normalmente en los
  \omterm{def:b3:topology:compact}{compactos} de $\C\setminus\Z$,
        \[
        \frac{\pi^2}{\sin^2(\pi z)} =
        \sum_{n\in\Z}\frac1{(z - n)^2} .
        \]
  \item Evalúese en $z = \frac12$:
        $\sum_{m\geq0}\frac1{(2m+1)^2} = \frac{\pi^2}8$;
        y recupérese $\zeta(2)$ una vez más separando los enteros por
        paridad. \item Verifíquese la identidad de duplicación
        $\tan\theta =
        \cot\theta - 2\cot(2\theta)$ y dedúzcase
        \[
        \pi\tan(\pi z) =
        \sum_{m\geq0}\frac{8z}{(2m+1)^2 - 4z^2} ,
        \]
        normalmente en los \omterm{def:b3:topology:compact}{compactos}
        que evitan $\frac12 + \Z$. \item Desarróllese alrededor de $0$
        ($\abs z < \frac12$; Fubini de nuevo): con $\lambda(s) =
        \sum_{m\geq0}(2m+1)^{-s}$,
        \[
        \pi\tan(\pi z) =
        \sum_{k\geq0}8\cdot4^k\,\lambda(2k+2)\,z^{2k+1} ;
        \]
        compárese con $\tan u = u + \frac{u^3}3 + O(u^5)$ para recuperar
        $\lambda(2) = \frac{\pi^2}8$ y obtener
        $\lambda(4) = \frac{\pi^4}{96}$, y contrástese después
        $\zeta(4) = \frac{\pi^4}{90}$ vía $\lambda(4) = (1 -
        2^{-4})\,\zeta(4)$. \item Verifíquese
        $\frac1{\sin\theta} = \cot\frac\theta2 -
        \cot\theta$ y dedúzcase
        \[
        \frac{\pi}{\sin(\pi z)} = \frac1z +
        \sum_{n\geq1}(-1)^n\,\frac{2z}{z^2 - n^2} .
        \]
        Compruébense los signos frente a los
        \omterm{def:b3:residues:singularities}{residuos} de
        $\pi/\sin(\pi z)$ en los enteros. \item Léase el coeficiente de
        $z$: $\eta(2) =
        \sum_{n\geq1}\frac{(-1)^{n-1}}{n^2} =
        \frac{\pi^2}{12}$, y confírmese la consistencia
        $\eta(2) = (1 - 2^{1-2})\,\zeta(2)$.
  \item (Final) Evalúese en $z =
        \frac12$ el desarrollo de la pregunta 20 y dedúzcase la
  \emph{fórmula de Leibniz}
        \[
        \frac\pi4 = 1 - \frac13 + \frac15 - \frac17 +
        \cdots
        \]
        Ciérrese con un breve párrafo: un núcleo por aritmética --- qué
        núcleo pone precio a qué familia de series y por qué todos ellos
        son estructuralmente ciegos a $\zeta(3)$.

\end{enumerate}
\medskip
\noindent\textbf{Parte VII --- La lista de precios \omterm{def:b3:complete:complete}{completa}: los números
de Bernoulli.}
\begin{enumerate}[resume]
  \item Combínese el desarrollo en fracciones simples de $\pi
        z\cot(\pi z)$ con la función generatriz de los números de
  Bernoulli ($\frac{w}{\eu^w - 1} =
        \sum_n\frac{B_n}{n!}w^n$, el~\cref{pb:b3:holomorphic:1}, Parte
  VI): a partir de
        \[
        \pi z\cot(\pi z) = \iu\pi z +
        \frac{2\iu\pi z}{\eu^{2\iu\pi z} - 1}
        \]
        (demuéstrese primero esta identidad), dedúzcase la forma cerrada
        \[
        \zeta(2k) = (-1)^{k+1}\,
        \frac{(2\pi)^{2k}\,B_{2k}}{2\,(2k)!}
        \qquad (k \geq 1).
        \]
  \item Verifíquese la fórmula frente a $B_2 = \frac16$, $B_4 =
        -\frac1{30}$, $B_6 = \frac1{42}$: recupérense $\zeta(2) =
        \frac{\pi^2}6$ y $\zeta(4) = \frac{\pi^4}{90}$, y calcúlese
  $\zeta(6) = \frac{\pi^6}{945}$. \item (Recursión de Euler)
  Desarróllense ambos miembros de
        $\bigl(z\cot z\bigr)' = \cot z - z(1 + \cot^2z)$ ---
        --- o elévese al cuadrado directamente la serie de la cotangente
        --- para demostrar
        \[
        \Bigl(k + \frac12\Bigr)\zeta(2k) =
        \sum_{j=1}^{k-1}\zeta(2j)\,\zeta(2k - 2j)
        \qquad (k \geq 2),
        \]
        y compruébese que calcula $\zeta(4)$ a partir de $\zeta(2)$ y
        $\zeta(6)$ a partir de $\zeta(2), \zeta(4)$ --- todos los
        valores pares de zeta a partir de la única semilla
        $\frac{\pi^2}6$, sin ninguna integración nueva.
\end{enumerate}
\end{problem}
